% Part: incompleteness
% Chapter: theories-computability
% Section: q-is-ce

\documentclass[../../../include/open-logic-section]{subfiles}

\begin{document}

\olfileid{inc}{tcp}{qce}

\olsection{$\Th{Q}$ تامة بين المجموعات \printtoken{S}{c.e.}}


\begin{thm}
$\Th{Q}$ !!{c.e.} لكنها غير قابلة للقرار. بل هي مجموعة !!{c.e.}
تامة.
\end{thm}

\begin{proof}
ليس من الصعب أن نرى أن $\Th{Q}$ !!{c.e.}، لأنها مجموعة
(شفرات) الجمل~$y$ التي يوجد لها برهان~$x$ في $\Th{Q}$:
\[
Q = \Setabs{y}{\lexists[x][\Prf[\Th{Q}](x,y)]}.
\]
لكننا نعلم أن $\Prf[\Th{Q}](x,y)$ قابلة للحساب (بل عودية بدائية)،
وكل مجموعة يمكن كتابتها بالصورة أعلاه قابلة للتعداد حسابيًا.

والقول إنها مجموعة تامة قابلة للتعداد حسابيًا يكافئ القول إن
$K\leq_m Q$، حيث $K=\Setabs{x}{!A_x(x)\downarrow}$. فلنبيّن إذن أن
$K$ قابلة للاختزال إلى $\Th{Q}$. لما كان محمول كليني $T(e,x,s)$
عوديًا بدائيًا، فهو قابل للتمثيل في $\Th{Q}$، ولتكن الصيغة الممثلة
$!A_T$. عندئذ، لكل $x$،
\begin{align*}
x \in K & \lif \lexists[s][T(x,x,s)] \\
& \lif \lexists[s][(\Th{Q} \Proves !A_T(\num x,\num x, \num s))]
  \\
& \lif \Th{Q} \Proves \lexists[s][!A_T(\num x, \num x, s)].
\end{align*}
وفي الاتجاه العكسي، إذا كانت
$\Th{Q}\Proves\lexists[s][!A_T(\num x,\num x,s)]$، فلا بد في
الواقع أن تكون الصيغة $!A_T(\num x,\num x,\num n)$ صادقة لعدد
طبيعي ما~$n$. فلو كذبت $T(x,x,n)$، لأثبتت $\Th{Q}$ الصيغة
$\lnot !A_T(\num x,\num x,\num n)$ لأن $!A_T$ تمثل~$T$. لكن
$\Th{Q}$ كانت ستثبت عندئذ صيغة كاذبة، وهذا تناقض. فلا بد أن تصدق
$T(x,x,n)$، ومن ذلك ينتج $!A_x(x)\downarrow$.

وخلاصة القول إن $x$ ينتمي إلى $K$، لكل $x$، إذا وفقط إذا أثبتت
$\Th{Q}$ الصيغة $\lexists[s][!A_T(\num x,\num x,s)]$. ومن ثم تكون
الدالة $f$ التي ترسل $x$ إلى (شفرة) الجملة
$\lexists[s][!A_T(\num x,\num x,s)]$ اختزالًا لـ$K$ إلى
$\Th{Q}$.
\end{proof}

\end{document}
